% theme: similarity_ratios_and_area_ratios
% 解答生成済み。表は三角形の平行線、裏は台形・平行四辺形の面積比。
\MajorQuestionLesson{相似比と面積比}
\SetRowSpace{2mm}

\TypeQuestion{$\Tri{ABC}$ の辺 $\Lseg{AB},\Lseg{AC}$ 上に点 $\mathrm{D},\mathrm{E}$ をとり、$\Lseg{DE}\Parallel\Lseg{BC}$ とする。$\Lseg{AD}:\Lseg{DB}=2:3,\ \Lseg{BC}=15\mathrm{cm}$ で、$\Tri{ABC}$ の面積は $100\mathrm{cm}^2$ である。}{similarity_area_ratio_triangle_cut}
\FigProblemBlock{50mm}{fig/lesson_10_triangle_cut}{
\QQ{線分 $\Lseg{DE}$ の長さを求めなさい。}{$6\mathrm{cm}$}
\QQ{$\Tri{ADE}$ の面積を求めなさい。}{$16\mathrm{cm}^2$}
\QQ{台形 $\mathrm{DBCE}$ の面積を求めなさい。}{$84\mathrm{cm}^2$}
\QQ{$\Tri{ADE}$ と台形 $\mathrm{DBCE}$ の面積比を求めなさい。}{$4:21$}
}{
}

\TypeQuestion{$\Tri{ABC}$ の辺 $\Lseg{AB},\Lseg{AC}$ 上に点 $\mathrm{D},\mathrm{E},\mathrm{F},\mathrm{G}$ をとり、$\Lseg{DE}\Parallel\Lseg{FG}\Parallel\Lseg{BC}$ とする。$\Lseg{AD}:\Lseg{DF}:\Lseg{FB}=2:1:3$ で、$\Tri{ABC}$ の面積は $180\mathrm{cm}^2$ である。}{similarity_area_ratio_multiple_bands}
\FigProblemBlock{54mm}{fig/lesson_12_parallel_bands}{
\QQ{$\Tri{ADE}$ の面積を求めなさい。}{$20\mathrm{cm}^2$}
\QQ{面積比$\Tri{ADE}:\Tri{AFG}$を求めなさい。}{$4:9$}
\QQ{四角形 $\mathrm{DFGE}$ の面積を求めなさい。}{$25\mathrm{cm}^2$}
\QQ{四角形 $\mathrm{FBCG}$ の面積を求めなさい。}{$135\mathrm{cm}^2$}
\QQ{$\Tri{ADE}$, 台形$\mathrm{DFGE}$, 台形$\mathrm{FBCG}$の面積比を求めなさい。}{$4:5:27$}
}{
}

\TypeQuestion{$\Tri{ABC}$ の辺 $\Lseg{AB},\Lseg{AC}$ 上に点 $\mathrm{D},\mathrm{E},\mathrm{F},\mathrm{G}$ をとり、$\Lseg{DE}\Parallel\Lseg{FG}\Parallel\Lseg{BC}$ とする。$\Lseg{AD}:\Lseg{DF}:\Lseg{FB}=1:2:3$ で、$\Tri{ABC}$ の面積は $216\mathrm{cm}^2$ である。$\Lseg{BG}$ と $\Lseg{CF}$ の交点を $\mathrm{O}$ とする。}{similarity_area_ratio_bands_butterfly}
\FigProblemBlock{58mm}{fig/lesson_12_parallel_bands_butterfly}{
\QQ{$\Tri{ADE}$ の面積を求めなさい。}{$6\mathrm{cm}^2$}
\QQ{面積比$\Tri{ADE}:\Tri{AFG}$を求めなさい。}{$1:9$}
\QQ{四角形 $\mathrm{DFGE}$ の面積を求めなさい。}{$48\mathrm{cm}^2$}
\QQ{四角形 $\mathrm{FBCG}$ の面積を求めなさい。}{$162\mathrm{cm}^2$}
\QQ{面積比$\Tri{FOG}:\Tri{COB}$を求めなさい。}{$1:4$}
\QQ{$\Tri{FBC}$ の面積を求めなさい。}{$108\mathrm{cm}^2$}
\QQ{$\Tri{FOG}$ の面積を求めなさい。}{$18\mathrm{cm}^2$}
}{
}

\AnswerColumnBreak
\LessonQuestionBreak

\Needspace{88mm}
\TypeQuestion{$\Lseg{AB}\Parallel\Lseg{CD}$ である台形 $\mathrm{ABCD}$ の対角線 $\Lseg{AC},\Lseg{BD}$ の交点を $\mathrm{O}$ とする。$\Lseg{AB}=6\mathrm{cm},\ \Lseg{CD}=15\mathrm{cm}$ で、$\Tri{AOB}$ の面積は $12\mathrm{cm}^2$ である。}{similarity_area_ratio_trapezoid}
\FigProblemBlock{40mm}{fig/lesson_10_trapezoid_area}{
\QQ{$\Tri{AOB}$ と $\Tri{COD}$ の相似比と面積比を、それぞれ求めなさい。}{相似比 $2:5$、面積比 $4:25$}
\QQ{$\Tri{COD}$ の面積を求めなさい。}{$75\mathrm{cm}^2$}
\QQ{$\Tri{BOC}$ の面積を求めなさい。}{$30\mathrm{cm}^2$}
\QQ{台形 $\mathrm{ABCD}$ の面積を求めなさい。}{$147\mathrm{cm}^2$}
}{
}

\Needspace{90mm}
\TypeQuestion{平行四辺形 $\mathrm{ABCD}$ の辺 $\Lseg{BC}$ 上に点 $\mathrm{E}$ をとり、$\Lseg{BE}:\Lseg{EC}=3:1$ とする。線分 $\Lseg{AE}$ と対角線 $\Lseg{BD}$ の交点を $\mathrm{P}$、対角線 $\Lseg{AC},\Lseg{BD}$ の交点を $\mathrm{Q}$ とする。平行四辺形 $\mathrm{ABCD}$ の面積は $112\mathrm{cm}^2$ である。}{similarity_area_ratio_parallelogram_inside}
\FigProblemBlock{56mm}{fig/lesson_09_parallelogram_04}{
\QQ{$\Tri{ABC}$ の面積を求めなさい。}{$56\mathrm{cm}^2$}
\QQ{$\Tri{ABE}$ の面積を求めなさい。}{$42\mathrm{cm}^2$}
\QQ{$\Tri{ABP}$ の面積を求めなさい。}{$24\mathrm{cm}^2$}
\QQ{$\Tri{APQ}$ の面積を求めなさい。}{$4\mathrm{cm}^2$}
}{
}

\Needspace{92mm}
\TypeQuestion{平行四辺形 $\mathrm{ABCD}$ で、点 $\mathrm{E}$ は辺 $\Lseg{DC}$ を点 $\mathrm{D}$ の側に延長した直線上にあり、$\Lseg{DE}:\Lseg{DC}=1:2$ である。線分 $\Lseg{BE}$ と対角線 $\Lseg{AC}$ の交点を $\mathrm{P}$、線分 $\Lseg{BE}$ と辺 $\Lseg{AD}$ の交点を $\mathrm{Q}$ とする。平行四辺形 $\mathrm{ABCD}$ の面積は $90\mathrm{cm}^2$ である。}{similarity_area_ratio_parallelogram_extension}
\FigProblemBlock{56mm}{fig/lesson_09_parallelogram_02}{
\QQ{$\Tri{ABC}$ の面積を求めなさい。}{$45\mathrm{cm}^2$}
\QQ{$\Tri{PBC}$ の面積を求めなさい。}{$27\mathrm{cm}^2$}
\QQ{$\Tri{APQ}$ の面積を求めなさい。}{$12\mathrm{cm}^2$}
\QQ{四角形$\mathrm{CPQD}$ の面積を求めなさい。}{$33\mathrm{cm}^2$}
}{
}
